\begin{abstract}
Multi-task model merging via unsupervised test-time adaptation (TTA) offers a flexible paradigm for consolidating specialized expert vectors on the fly. However, we expose a critical vulnerability in unconstrained layer-wise merging optimizers: they suffer from what we term the \emph{Overfitting-Optimizer Paradox} (an overfitting phenomenon under test-time adaptation), where high degrees of freedom allow the model to minimize surrogate test-time entropy by memorizing transient, high-frequency transductive noise, leading to severe representation collapse. To resolve this, we propose \textbf{ChebyMerge}, a mathematically rigorous framework that projects the layer-wise coefficient space onto an orthogonal, low-dimensional subspace spanned by Chebyshev polynomials of the first kind. By framing merging as a continuous spectral approximation problem, ChebyMerge achieves minimax-optimal uniform approximation, concentrates representation resolution at the highly sensitive model boundaries, and filters out transductive high-frequency noise. Mathematically, ChebyMerge yields an exceptionally well-conditioned optimization landscape, improving the Gram matrix condition number by up to 3,527$\times$ over standard monomial power bases (e.g., PolyMerge) for cubic parameterizations. Using a physically-grounded coupled non-convex loss-landscape simulator designed to emulate the sensitivities, functional couplings, and transductive noise of deep vision-language networks, we demonstrate that ChebyMerge stabilizes on-the-fly optimization, prevents representation collapse, and achieves robust, predictable generalization. Finally, we expose and analyze the \emph{Conditioning-Generalization Paradox} (representing the classic tension between numerical conditioning and regularization), demonstrating that while monomial ill-conditioning can act as an accidental, uncontrolled regularizer (spectral damping), ChebyMerge is a mathematically superior and more controllable framework that successfully decouples numerical stability from parameter regularization.
\end{abstract}
